%% ============================================================
%% sections/s2-findings.tex
%% H. R. 9510 Bill v3.0 - SEC. 2. FINDINGS; THE EVIDENCE TO LAW RECORD.
%% Rendered visuals: Figure 1 (four-work evidence-to-law lineage), Table 1 (the
%% four works at a glance), Figure 2 (accelerated timeline). The fourteen findings
%% of Bill v2.0 are condensed into a lettered narrative carried by the three
%% visuals, in place of a long numbered list.
%% ============================================================
\billsec{SEC. 2.}{FINDINGS; THE EVIDENCE TO LAW RECORD.}\phantomsection\label{toc:s2}

\uscprov{1}{\textbf{(a) In General.} Congress finds that an automated process
must verify robot-patient interaction code before that code is generated or
executed in a Physical AI oncology clinical investigation, and that the record of
the VVUQ-01 through VVUQ-04 developments, visualized in this section, establishes
both the need for and the feasibility of that requirement \cite{kawchak2026national,
vvuq05-figures-bill}. The lineage in Figure 1 shows how engineering evidence
became proposed law and then a precise amendment.}

\begin{asciifigsmall}
  VVUQ-01                  VVUQ-02                  VVUQ-03                  VVUQ-04
  (method + pipeline)      (hard engineering)       (proposed law)           (precise amendment)
  +-------------------+    +-------------------+    +-------------------+    +-------------------+
  | inputs/ ingest    |    | instructions/     |    | template-paper/   |    | instruct-bill/    |
  | execution/ 51 tst | -> | codegen/ 10 gates | -> | 21 CFR Part 312   | -> | 10 md + 5 bib     |
  | image-instruct/   |    | execution/ 172tst |    | instruct-paper/   |    | template-bill/    |
  | imagegen/ 10 figs |    | imagegen/ 15 figs |    | draft-paper/      |    | draft-bill/       |
  | draft -> full ->  |    | draft -> full ->  |    | full-paper/       |    | full-bill/        |
  | final-paper       |    | final-paper       |    | final-paper(bill) |    | final-bill        |
  +-------------------+    +-------------------+    +-------------------+    +-------------------+
            | DOI                    | DOI                    | DOI                    | DOI
            v                        v                        v                        v
       zenodo.20372501          zenodo.20421754          zenodo.20454870          zenodo.20485580
            \________________________\________________________\_______________________/
                                                 |
              evidence flows forward: each work cites the records of the prior
              works (VVUQ-02 cites VVUQ-01; VVUQ-03 cites VVUQ-01, VVUQ-02, and
              the national platform; VVUQ-04 cites VVUQ-03 and the Title 21 Code)
\end{asciifigsmall}
\figcaption{Figure 1. Four-work evidence-to-law lineage. The method and pipeline
(VVUQ-01), the hard engineering proof (VVUQ-02), the proposed bill (VVUQ-03), and
the precise amendment (VVUQ-04), each released to a Zenodo DOI.}

\uscprov{1}{\textbf{(b) The Embodiment Problem and the Sequencing Gap.} Physical
AI oncology investigations place autonomous and semi-autonomous robots in direct
physical contact with patients, and the behavior of a robot at the moment of
contact is determined by software that is increasingly machine-generated. In a
disembodied pipeline an error yields a wrong number a human can catch downstream;
in an embodied surgical system the same error class yields a wrong motion a human
cannot catch in time. No Federal law requires the verification of robot-patient
interaction code to clear before that code is generated or executed, because the
order of operations is not regulated \cite{kawchak2026vvuq01, kawchak2026national}.
Table 1 records what each of the four works produced.}

{\footnotesize
\begin{xltabular}{\textwidth}{@{}L{1.2cm} Y L{1.6cm} Y L{1.7cm}@{}}
\hline
\textbf{Work} & \textbf{What it is} & \textbf{Releases} & \textbf{Headline
assurance result} & \textbf{End-goal DOI}\\
\hline
\endfirsthead
\hline
\textbf{Work} & \textbf{What it is} & \textbf{Releases} & \textbf{Headline
assurance result} & \textbf{End-goal DOI}\\
\hline
\endhead
\hline
\endlastfoot
VVUQ-01 & Method paper plus the working pipeline and an execution record that hold
VVUQ higher than code generation & v0.1.0 to v0.6.0 & 51 of 51 tests; a 1 ACCEPT,
5 BLOCK, 1 ESCALATE gate surface; 2.5x acceleration &
\href{https://doi.org/10.5281/zenodo.20372501}{20372501}\\
\hline
VVUQ-02 & A generated ten-gate humanoid surgical codebase, its autonomous
execution record, and the manuscript built from them & v0.7.0 to v1.2.0 & 172 of
172 tests; 10 gates; 32 of 32 sweep iterations clear; composite mean 93.56 &
\href{https://doi.org/10.5281/zenodo.20421754}{20421754}\\
\hline
VVUQ-03 & The standalone Physical AI Oncology Trial Bill (H. R. 9510),
synthesized from the recorded evidence & v1.3.0 to v1.4.0 & 15 sections; 21
body-width tables; 60 references; verification before generation as statute &
\href{https://doi.org/10.5281/zenodo.20454870}{20454870}\\
\hline
VVUQ-04 & The bill recast as a precise amendment to the Federal Food, Drug, and
Cosmetic Act, current through Public Law 119-93 & v2.0.0 to v2.3.0 & New section
360e-5 plus ten conforming amendments; legal crosswalk current to May 31, 2026 &
\href{https://doi.org/10.5281/zenodo.20485580}{20485580}\\
\hline
\end{xltabular}}
\figcaption{Table 1. The four works at a glance, with each work's role, release
range, headline assurance result, and end-goal DOI.}

\uscprov{1}{\textbf{(c) Feasibility Is Demonstrated.} Automated verification of
robot-patient interaction code is feasible today. A pipeline held the
verification fraction at 1.0, accepting one candidate artifact and blocking five,
and a ten-gate assurance suite over control software for a surgical humanoid
cleared reproducibly from a fixed seed, each gate bound to a published external
standard \cite{kawchak2026vvuq01, kawchak2026vvuq02}.}

\uscprov{1}{\textbf{(d) The Closest Existing Analogues.} The predetermined change
control plan authority in section 515C of the Federal Food, Drug, and Cosmetic
Act (21 U.S.C. 360e-4) and the Food and Drug Administration final guidance on such
plans for artificial-intelligence-enabled device software functions are the
closest existing analogues to validating a change before deployment, but each
remains voluntary, submission-based, and neither automated nor required in
advance of generation. The Quality Management System Regulation (part 820 of
title 21, Code of Federal Regulations), effective February 2, 2026, and the
electronic-records requirements of part 11 of such title supply quality and
record-integrity scaffolding but do not impose a rule that verification precede
generation \cite{fda-pccp-2024, fr-qmsr-2024}.}

\uscprov{1}{\textbf{(e) A Recognized Standards Basis.} Recognized consensus
standards, including ASME V\&V 40-2018, the Food and Drug Administration guidance
assessing the credibility of computational modeling, NASA-STD-7009A, IEEE Std
7009-2024, IEC 80601-2-77, ISO 14971, and ISO/TS 15066, supply the credibility
framework, dispersion bounds, fail-safe requirements, and biomechanical contact
limits to which an automated gate may be bound \cite{asmevv40, fda-credibility-2023,
nasastd7009, isots15066}.}

\uscprov{1}{\textbf{(f) Human-Over-AI Review Is Already Widely Adopted.}
Human-over-AI review is already widely adopted in State law, including the
California Physicians Make Decisions Act (Senate Bill 1120), the Illinois Wellness
and Oversight for Psychological Resources Act (House Bill 1806), the Texas
Responsible Artificial Intelligence Governance Act (House Bill 149), and Maryland
House Bill 820, and the strongest Federal precedent for mandatory human-over-AI
oversight is the Medicare Advantage artificial-intelligence guardrail (section
422.101(c) of title 42, Code of Federal Regulations) \cite{ca-sb1120, il-hb1806,
tx-hb149, md-hb820, cms-4201-ma}.}

\uscprov{1}{\textbf{(g) Modernization, Nondiscrimination, and Classification.} As
clinical investigation moves toward real-time and decentralized designs governed
by harmonized good clinical practice, manual review cannot keep pace with the
volume of machine-generated control code, and an automated gate that clears
before generation serves that modernization. The acceleration the four works show
is set against the conventional path in Figure 2. Federal nondiscrimination law,
including section 92.210 of title 45, Code of Federal Regulations, title VI of the
Civil Rights Act of 1964 (42 U.S.C. 2000d), and the Americans with Disabilities
Act of 1990 (42 U.S.C. 12101 et seq.), requires that a patient-facing algorithm
minimize the risk of bias and respect accessibility, and the
artificial-intelligence taxonomy maintained by the American Medical Association
provides a graduated spine for matching verification rigor to the degree of
autonomy of a device \cite{iche6r3, cfr-92-210, usc-titlevi, usc-ada,
cpt-appendix-s}.}

\begin{asciifig}
  Conventional path (months to years)        This directory (about 11 days)
  +-----------------------------------+      +-------------------------------------+
  | Method paper drafting   ~months   |      | VVUQ-01  v0.1.0 -> v0.6.0  ~5 days  |
  | Verified robotics code  6-18 mo   |  vs  | VVUQ-02  v0.7.0 -> v1.2.0  ~2 days  |
  | Peer-style manuscript   ~months   |      | VVUQ-03  v1.3.0 -> v1.4.0  ~1 day   |
  | Federal bill + counsel  mo-years  |      | VVUQ-04  v2.0.0 -> v2.3.0  ~2 days  |
  +-----------------------------------+      +-------------------------------------+
  years of serial effort                     days of autonomous, real-time work
                                             reaching the final VVUQ-04 amendment
\end{asciifig}
\figcaption{Figure 2. Accelerated timeline versus conventional methods. The
conventional months-to-years path beside the about eleven days of autonomous,
real-time work that produced VVUQ-01 through VVUQ-04.}
